% ===== 原文 pp.100–102（PDF p.110–112） =====
\chapter{紧化在亏格 \texorpdfstring{$g\ge 2$}{g>=2} 黎曼面上的 Vafa-Witten 理论}

% [原文 p.100]
此前我们讨论过 Vafa-Witten 理论在圆上紧化到三维。现在我们来讨论 Vafa-Witten 理论在黎曼面 $C$ 上紧化到二维。约化的细节强烈依赖于 $C$ 的亏格 $g_C$，原因在于 Hitchin 方程解的模空间的虚拟维数在 $g_C\ge 2$ 时为正、在 $g_C=1$ 时为零、在 $g_C=0$ 时为负；这大致对应于如下事实：当 $g_C\ge 2$ 时，理论低作用量方程的一般解会破缺全部规范对称性，而 $g_C=0,1$ 时则不然。为避免未破缺规范对称性带来的复杂性，本章处理 $g_C\ge 2$ 的情形，并把分析限制在不可约构形上。（这相当于切除二维有效理论靶空间中的奇异点。）

我们的分析将与文献 [1] 中对 GL-扭转理论作二维约化时采用的分析相当类似。在那里，$t=\ii$ 时的有效二维 TFT 被确定为 B-模型，$t=1$ 时的有效二维 TFT 被确定为 A-模型。然而我们将会发现，Vafa-Witten 理论的约化给出一个新的二维 TFT：第 7 章的修正 A-模型。（这个理论在 §6.4 分析面算子时也曾短暂出现。）“Vafa-Witten 理论在黎曼面上约化为以 $\cM_H(G,C)$ 为靶空间的修正 A-模型”这一论断最早是在文献 [20] 的附录 B 中得到的。
% [Tutor Note: p.100 关于 $g_C$ 的讨论可用 p.104 的 $\dim\cM_H(G,C)=4(g_C-1)\dim G$ 核对：$g_C=1$ 给 0，$g_C=0$ 给负值（虚拟维数），与原文说法一致；另注意此处“低作用量”而非“低能量”，因为工作在欧氏号差下（见 p.101 脚注 1）]

\section{约化 Vafa-Witten 作用量}

考虑定义在 $\Sigma\times C$ 上的 Vafa-Witten 理论，其中 $\Sigma$ 与 $C$ 都是闭黎曼面，亏格分别为 $g_\Sigma$ 与 $g_C$，并假定 $g_C\ge 2$。我们希望求出 $C$ 的尺寸收缩为零的极限下的有效理论。更确切地说，我们希望求出 $\Sigma$ 上的一个二维 TFT，并给出 $\Sigma$ 上的二维场构形与 $\Sigma\times C$ 上低作用量四维场构形之间的一个字典，使得二维与四维作用量近似相等。\footnote{由于我们在欧氏号差下工作，有效理论描述的是低作用量的构形，而非低能量的构形。}

% [原文 p.101]
由 $C$ 上联络的不可约性假设，全部规范对称性都将被破缺，二维 TFT 将简单地是映射 $\Sigma\to\cM_H(G,C)$ 的（非规范）$\sigma$-模型。

把 $C$ 上的度规 $h_C$ 用小参数 $\epsilon$ 重新标度
\[
h_C\to h'_C=\epsilon\,h_C
\]
并考虑 $\epsilon\to 0$ 的极限，此时 $C$ 的体积与 $\Sigma$ 的体积相比很小。用 $\Sigma$ 上的实坐标 $\sigma^1,\sigma^2$ 与 $C$ 上的实坐标 $\sigma^3,\sigma^4$（按 §7.1 取与 $\Sigma$、$C$ 上复结构相适配的坐标），重新标度后的度规取形式
\[
h_\Sigma\oplus h_C=\diag\left(\ee^{2f_\Sigma},\ee^{2f_\Sigma},\epsilon\,\ee^{2f_C},\epsilon\,\ee^{2f_C}\right)
\]
其中实函数 $f_\Sigma(\sigma^1,\sigma^2)$ 与 $f_C(\sigma^3,\sigma^4)$ 是共形因子。

回忆第 3 章：Vafa-Witten 理论包含玻色型 0-形式 $\varphi$ 与 $\sigma$，以及自对偶 2-形式 $B$ 与联络 $A$。场 $B_{12}$ 关于 $\Sigma$ 与 $C$ 上的旋转都是 0-形式。场 $B_{13}$ 与 $B_{14}$ 按 $C$ 上 1-形式的分量变换，同时也——这一点至关重要——按 $\Sigma$ 上 1-形式的分量变换。分量 $B_{34}$、$B_{24}$、$B_{23}$ 由自对偶性用 $B_{12}$、$B_{13}$、$B_{14}$ 确定；为叙述简便，我们忽略费米子。
% [Tutor Note: p.101 “$B_{34},B_{24},B_{23}$ 由自对偶性确定”处省略了指标升降：欧氏号差下取 $\epsilon_{1234}=+1$ 且 $(*B)_{ij}=\frac12\epsilon_{ij}{}^{kl}B_{kl}$，自对偶 $*B=B$ 给出 $B_{34}=B^{12}$、$B_{24}=-B^{13}$、$B_{23}=B^{14}$。原文写 $B_{34},B_{24},B_{23}$ 而右边用 $B_{12},B_{13},B_{14}$，严格说应带升指标（$h$ 非平凡时二者不等），这是符号层面的简写]

在重新标度之后，Vafa-Witten 作用量 (3.9) 中的不同项会乘上 $\epsilon$ 的不同幂次。在 $\epsilon\to 0$ 极限下，低作用量场构形就是那些使乘以最负幂次 $\epsilon$ 的项为零的构形。这样便得到如下低 $\epsilon$ 方程：
\begin{gather}
\left(\ee^{2f_\Sigma}\right)F_{34}-[B_{13},B_{14}]=0 \notag\\
D_3\left(\ee^{-2f_C}B_{13}\right)+D_4\left(\ee^{-2f_C}B_{14}\right)=0 \tag{8.1}\\
D_3\left(\ee^{-2f_C}B_{14}\right)-D_4\left(\ee^{-2f_C}B_{13}\right)=0 \notag
\end{gather}
% [Tutor Note: 式 (8.1) 的 $\epsilon$ 幂次计数是原文省略的论证：$h_C$ 带因子 $\epsilon$ 后，每个 $C$ 方向逆度规贡献 $\epsilon^{-1}$，而积分测度 $\sqrt{h}$ 贡献 $\epsilon^{+1}$；故含两个 $C$ 指标的项（$F_{34}^2$、$D_3(\cdot)^2$ 型）整体给 $\epsilon^{-1}$，取 $\epsilon\to0$ 令其系数为零即得 (8.1)。自洽性检验：取 $f_\Sigma=f_C=0$、$\epsilon=1$，第一式化为 $F_{34}=[B_{13},B_{14}]$，与 p.102 的 $F_C=\frac12[\varphi,\varphi]$ 一致]

此外，标量 $B_{12}$、$\varphi$ 与 $\sigma$ 被要求沿 $C$ 协变常值；例如
\[
D_3B_{12}=D_4B_{12}=0 .
\]

0-形式的协变常值性可解释为：它们生成保持场 $A_C$ 不动的无穷小规范变换。但按 $A_C$ 的构形破缺全部规范对称性的假设，这又要求协变常值的 0-形式必须为零。因此，场 $\sigma$、$\varphi$ 与 $B_{12}$ 不再作为有效理论的动力学场存活下来。
% [Tutor Note: p.101 末句可这样补出：设 $b$ 是任一 0-形式（如 $\sigma,\varphi,B_{12}$），若沿 $C$ 有 $D_C b=0$，则由 $b$ 生成的无穷小规范变换使 $A_C$ 不变（$\delta A_C=D_C b=0$），即 $b$ 属于 $A_C$ 的稳定化子的李代数。不可约性假设恰是“稳定化子最小”，故此时只能 $b=0$；这正是原文“它们不存活下来”的依据。可对照 p.114 的 $\delta\mu_I(\epsilon)=-\iota_{V(\epsilon)}\omega_I$ 理解“0-形式生成规范变换”的含义]

% [原文 p.102]
至于场 $A_3$、$A_4$、$B_{13}$、$B_{14}$，上述方程可以写成如下形式
\begin{gather}
F_{34}-[\varphi_3,\varphi_4]=0 \notag\\
D_3\varphi^3+D_4\varphi^4=0 \tag{8.2}\\
D_3\varphi_4-D_4\varphi_3=0 \notag
\end{gather}
其中 $\varphi$ 是 $C$ 上重新标度后的 1-形式，由下式给出
\[
\varphi=\ee^{-f_\Sigma}\left(B_{13}\,\dd\sigma^3+B_{14}\,\dd\sigma^4\right).
\]
% [Tutor Note: p.102 式 (8.2) 第二式的上标 $\varphi^3,\varphi^4$ 与第一、三式的下标 $\varphi_3,\varphi_4$ 不一致，属原文笔误/记法不统一：$\varphi\in\Omega^1(C,\mathrm{ad}\,P_C)$ 是 1-形式，其分量应为下标（与 (8.1) 中 $B_{13},B_{14}$ 的写法一致，且 $\varphi=\varphi_3\dd\sigma^3+\varphi_4\dd\sigma^4$）。判断依据：$D_C\varphi=0$ 与 $D_C^*\varphi=0$ 按分量写出应为一组同型的方程，而 $\varphi$ 作为 $C$ 上 1-形式的分量按约定写成下标（与 (8.1) 中 $B_{13},B_{14}$ 相同）；此处照抄原文，不擅改]

这些正是黎曼面 $C$ 上主 $G$-丛 $P_C$ 的联络 $A_C$ 与截面 $\varphi\in\Omega^1(C,\mathrm{ad}\,P_C)$ 所满足的方程
\[
F_C=\tfrac12[\varphi,\varphi]=D_C^*\varphi=D_C\varphi=0,
\]
它们已由 N. Hitchin [64] 研究过。这些方程的解模去规范变换所得的空间是一个有限维空间 $\cM_H(G,C)$，即 Hitchin 模空间。由于在施加上述方程之后 $A_C$ 与 $\varphi$ 对 $(\sigma^1,\sigma^2)$ 的依赖仍是任意的，这些场就定义了一个映射
\[
\Phi:\Sigma\to\cM_H(G,C).
\]
其实，这个说法并不完全正确，因为 $\varphi$ 不是 $\Sigma$ 上的标量：在 $\Sigma$ 的 $U(1)$ 结构群作用下，场 $\varphi_w$ 会被旋转成 $\varphi_{\bar w}$。因此，对 $(A_C,\varphi)$ 定义了一个截面
\[
\Phi\in\Omega^0\left(\Sigma,\;P_\Sigma\times_{U(1)}\cM_H(G,C)\right),
\]
它是与 $U(1)$ 标架丛 $P_\Sigma$ 相伴的纤维丛的截面，相伴所用的 $U(1)$ 作用是
\[
\cU_1:\cM_H(G,C)\to\cM_H(G,C),
\]
由下式给出
\[
\cU_1:(A_w,A_{\bar w},\varphi_w,\varphi_{\bar w})\to(A_w,A_{\bar w},\ee^{\ii\alpha}\varphi_w,\ee^{-\ii\alpha}\varphi_{\bar w})
\]
其中 $\ee^{\ii\alpha}\in U(1)$。

支配场 $A_\Sigma$、$A_C$ 与 $\varphi$ 的有效作用量将由 Vafa-Witten 作用量的 $\epsilon^{-1}$ 展开中的次领头项给出；也就是说，我们施加低 $\epsilon$ 方程并丢弃当 $\epsilon\to0$ 时变得可忽略的项。这样便得到如下截断作用量
\begin{equation}
\begin{aligned}
S_{\text{trunc.}}=-\frac{1}{e^2}\int_M \dd^4\sigma\,\Tr\Big\{&(F_{13}-F_{24})^2+(F_{14}+F_{23})^2\\
&+\left(\partial_1\varphi_3+(\partial_1f_\Sigma)\varphi_3+\partial_2\varphi_4+(\partial_2f_\Sigma)\varphi_4\right)^2\\
&+\left(\partial_1\varphi_4+(\partial_1f_\Sigma)\varphi_4-\partial_2\varphi_3-(\partial_2f_\Sigma)\varphi_3\right)^2\\
&+2F_{13}F_{24}-2F_{14}F_{23}\Big\}\Big|_{(A_3,A_4,\varphi_3,\varphi_4)\ \text{satisfy Hitchin equations on }C}
\end{aligned}
\tag{8.3}
\end{equation}
% [Tutor Note: p.102 的符号复用易混：p.101 的 $\varphi$ 是（被丢弃的）$0$-形式标量，而 (8.2) 起的 $\varphi$ 是重整标度后的 1-形式 $\varphi=e^{-f_\Sigma}(B_{13}\dd\sigma^3+B_{14}\dd\sigma^4)$；原文未提示这次重新定义，翻译时保留原符号以免与原文式号脱节。核对 (8.3) 项数：前两项是 $F$ 的 $\Sigma C$ 混合分量，后一项系数 2 与 $F_{13}F_{24}-F_{14}F_{23}$ 的拓扑（瞬子）来源对应]
